\newcommand{\pG}{\partial G}

\section{Applications.}
\subsection{Margin of Victory Bounds.}
The \emph{Margin of Victory} (MoV) of an election is the smallest number of ballots that must be modified in order to change the outcome of the election. 
Computing the MoV is useful for election reporting -- voters are bound to be curious how close the election was -- and also for election auditing (cf. \S\ref{sec:mismatch}).

STV is among the most difficult election rules to find a margin for. 
Even approximating the MoV is known to be NP-hard\footnote{Meaning complexity asymptotically greater than $O(N C \log(C))$ for some profile.} for STV; indeed, just deciding whether MoV$=1$ is NP-hard \cite{xiaComputingMarginVictory2012}.%\footnote{Note, in view of what follows in this section, that this implies audit graph construction is NP-hard -- any construction algorithm must cost $ \gg\!O(N m \log(m))$ operations for some profile.} in general.

Computing a lower bound on the MoV is more productive. 
There is already substantial work on this problem using a branch-and-bound Mixed-Integer Non-Linear Program (MINLP) \cite{blomAdvancesSTVMargin2026}, and such a bound may still be used to formulate an RLA for the election \cite{ekDoingMoreLess2026}.
Audit graphs offer an alternative jumping-off point to the computation of these bounds, which might be used in tandem with the MINLP approach to produce tighter bounds.
%In particular, they allow for an optimization of the branch-and-bound phase of the algorithm...

The approach follows the intuition of \S\ref{sec:plausible}: a coherent $M$-plausible audit graph $G$ is a witness that manipulating fewer than $M/2$ ballots can not change the outcome of the election, \emph{provided that} such a manipulation can not shift the tallies in any vertex of $G$ by an $L^1$ distance of more than $M$.
This provision holds as long as the tally function satisfies some notion of ``conservation of ballot mass:'' the sum of a single ballot's contribution to all tallies is at most $1$.

This is a principle that the STV tabulation algorithm was specifically encoded to preserve.
Ballots start the election with a weight of $1$, and transfer values are ordinarily computed to live in $[0,1]$.
The only difference in the context of audit graph construction is the risk of negative transfer values: when a winner's tally is contained in $[q-M, q)$, the edge corresponding to her seating will be considered plausible, even though she did not make quota.
In any vertices downstream from this seating, the winner's transfer value, computed as in Equation (\ref{eq:transfer}) of Appendix \ref{sec:appendix}, will be contained in $[-M/(q-M), 0)$.
This could break the conservation of ballot mass if the lower endpoint of this interval is less than $< -1$, prompting the following statement:
\begin{theorem}\label{thm:reasonable}
As long as $M \leq q/2$, a coherent $M$-plausible graph $G$ witnesses that the MoV of an election is greater than $M/2$.
\end{theorem}
\begin{proof}
	The bound comes from solving $-M/(q-M)=-1$; as long as it holds, all transfer values in $G$ are in $[-1,1]$, and modifying a single ballot can not shift the tallies by more than $2$ in $L^1$ distance. 

	If fewer than $M/2$ ballots are modified, the triangle inequality then implies the tallies in any vertex stay within a radius-$M$ $L^1$ ball of their recorded values in $G$, so the new natural path can not escape $G$.\qed
\end{proof}

It bears emphasis that the construction of plausible graphs is not optimized for MoV computation. % -- their purpose was first and foremost to generate RLAs.
In particular, the generosity described in \S\ref{sec:plausible} means that a plausible path contained in a graph $G$ might not be realizable by a perturbation of $\leq M/2$ ballots.
The path might need to change one set of $M/2$ ballots in one round of the election, and another disjoint set in the next.

A shrewder MoV-bounding approach would be to use an audit graph to detect the most plausible election path leading to an alternative winner set, then run the MINLP algorithm of \cite{blomAdvancesSTVMargin2026} on the prefix of that path leading to its edge with the largest plausibility threshold. 
Then, replace said threshold with the program's output, and repeat the process with the next-most plausible alternative path.

Imperfect as they are, plausible graphs can still be close to sharp in some cases.
For example, the plausible graph in Figure \ref{fig:vic} witnesses a lower bound $\text{MoV}\geq 11,500$ for its election, which closes the gap with a known upper bound of $\text{MoV}\leq 16,021$ computed using the ConcreteSTV algorithm of \cite{blomAdvancesSTVMargin2026}.

%A more efficient use of the audit graph construction for MoV-bounding could follow these lines:
%\begin{enumerate}
%    \item Increase $M$ until there is a 
%\end{enumerate}

\subsection{Mismatch-Based Audits.}\label{sec:mismatch}
The mismatch-based approach generates an RLA by verifying that fewer than a MoV worth of votes were reported incorrectly \cite{ekDoingMoreLess2026}. 
This is a so-called ``ballot-comparison'' audit: we have access to a digital \emph{Cast Vote Record} (CVR) indicating how each vote was interpreted by the tabulation, as well as a physical \emph{Manual Vote Record} (MVR) representing a paper trail which can be compared to the record by an Auditor.
The data input into the RLA test process is a sequence of CVR-MVR pairs which are compared for mismatches.

Audit graphs allow for a new subtlety in mismatch-audit design, by giving the auditing process a way of filtering out ``irrelevant'' mismatches. 
For example, in a single-seat contest with three candidates $\CCC = \{a, b, c\}$, if there is a separate test process already verifying that $c$ can not win the election, then a mismatch between a ballot reported as $(a,b)$ in the CVR but sampled as $(a, c)$ in the MVR can safely be ignored; either way, the ballot will count towards candidate $a$'s tally when it matters.

To formalize this, given a fixed plausible graph $G\subset \oO$, we call an edge $\ve = \ou\,\ov\in E(\oO)$ an \emph{escape} edge if $\ou\in G$ but $\ov\notin G$. 
We use the notation $\partial G$ to denote the set of all escape edges of $G$.

We define our \emph{global null hypothesis} as $H_{\star}=$ ``the natural election path under the MVR escapes the graph $G$'', and for each $\ve\in\pG$ we define a \emph{local null hypothesis} $H_{\ve}=$ ``the edge $\ve$ is natural.''
Under the SHANGRLA framework, a valid test of $H_{\star}$ at level $\alpha$ is to reject each of the $H_{\ve}$ at level $\alpha$, since $H_{\star}\subset \bigcup_{\ve\in\pG} H_{\ve}$. In particular, there is no need to adjust for multiplicity -- cf. \S4.1 of \cite{starkSetsHalfaverageNulls2020}.

To reject a local null $H_{\ve}$ where $\ve = \ou\, \ov$, it is always sufficient to verify that a given affine function of the tallies in $\ou$ which was recorded as positive in the CVR stays positive in the MVR. We call this affine function the \emph{critical margin} controlling the plausibility threshold of $\ve$ -- in practice, it can take one of three forms, depending on $\ve$:
\begin{itemize}
	\item The critical margin can be a \emph{candidate-to-candidate} margin, of the form $M_{cl} = T(c) - T(l)$, where $c, l \in \HHH(\ou)$ are two distinguished candidates.
	\item Alternatively, the critical margin can be a \emph{candidate-below-quota} or \emph{candidate-above-quota} margin, $M_{qc} = q-T(c)$ or $M_{cq} = T(c) -q$, where $c\in\HHH(\ou)$ is a distinguished candidate.
\end{itemize}
For example, if $\ve$ is a seating edge corresponding to winners $S\subset \HHH(\ou)$, a valid critical margin is the candidate-below-quota margin $M_{qc}$ corresponding to the candidate $c\in S$ with the lowest tally according to $v$. 
If $\ve$ is an elimination edge corresponding to $c\in \HHH(\ou)$, a valid critical margin might be the candidate-to-candidate margin $M_{cl}$ where $l\in \HHH(\ou)$ is the hopeful candidate with the lowest recorded tally in $\ou$, or alternatively the candidate-above-quota margin $M_{wq}$ where $w\in \HHH(\ou)$ is the candidate with the highest tally in excess of quota.

Depending on which of these margins we are working with, we can drastically reduce the dimension of the parameter space established in Equation (\ref{eq:parameters}) of Appendix \ref{sec:appendix} from $(|\HHH(\ou)|+1)2^{d(\ou)}$ to either $3\cdot 2^{d(\ou)}$ or $2^{d(\ou)+1}$.
For example, in a candidate-to-candidate margin $M_{cl}$, any parameter $t_{S;o}$ where $o\notin \{c,l\}$ has an identical impact, as long as $S$ is the same -- so they may all be treated as a single ``reduced parameter.''

Then at the level of a local null $H_{\ve}$, a CVR-MVR pair should only be considered a mismatch if their interpretation in this reduced parameter space differs.
To reject $H_{\ve}$, it is sufficient to check that there are not enough of these mismatches to reverse the critical margin computed by the CVR.
This in turn can be done via any of the standard test supermartingales integrated in the SHANGRLA framework, such as ALPHA or COBRA \cite{starkALPHAAuditThat2023c,spertusCOBRAComparisonOptimalBetting2023}.
We use the latter for the results in \S\ref{sec:results}.

The advantage of considering mismatches at the local rather than global level is a dilution of the overall noise level between the CVR and MVR.
Assuming the noise is not adversarial, most CVR-MVR discrepancies will only affect a fraction of the escape edges in $\pG$, leading to a more efficient test process.
\subsection{Delta-Method Audits.}\label{sec:delta}

We might also consider alternative test processes to reject the local null hypotheses $H_{\ve}$ of the previous section.
In particular, a test process which succesfully uses the information contained in each mismatch -- which candidate the mismatch helps or hurts -- is bound to be more efficient in the non-adversarial setting.

However, the non-linearity of the tally function in Equation (\ref{eq:tally}) of Appendix \ref{sec:appendix} means that the standard test supermartingales can't be applied off the shelf.
This is because the direct impact of a sampled mismatch on the critical margin depends not only the mismatch itself, but also the sampling history of other mismatches.
For example, sampling the same mismatch repeatedly can have a self-reinforcing impact on the margin.
We must adapt our test processes to these non-linear margins.

One such test uses the delta method to construct confidence intervals for each critical margin, and rejects the local null if this interval is positive. 
This involves bounding the variances and covariances of the sample means of each of the reduced parameters from \S\ref{sec:mismatch} and multiplying by an appropriate gradient matrix to obtain a variance for the critical margin, from which we build a confidence interval -- cf. \cite{heitzmannGraphBasedAuditsMeek2026a} for details.

This test process uses a standard Wald Confidence Interval, and as such it requires a fixed sample size -- it does not share the arbitrary stopping times of test supermartingales.
On the other hand, it is more resilient to noise than tests like mismatch-based audits, because in the regime of low noise it replaces degenerate sample variances with worst-case stochastic bounds which are more sensitive to the sample size than the (low) number of sampled mismatches.
But this also means that the Delta method performs worse for edges leaving higher-degree vertices, since these variance bounds are scaled by the number of parameters in Equation \ref{eq:parameters} to give the variance of the critical margin, and that number of parameters grows exponentially in the degree of the base vertex.

In summary, the delta method has some advantages and drawbacks compared to mismatch-based audits. 
It should not be taken as a final ``best practice'' test process to audit the assertions that audit graphs give rise to, but rather as a first pass attempt at using the information contained within a CVR-MVR comparison, which will be improved (indeed, already has been improved) in due time.
We show some results comparing the performance of the mismatch and delta methods in the next section. %, although better test processes will surely arise eventually
\subsection{Results.}\label{sec:results}
We constructed audit graphs and generated RLAs for a representative cross-section of STV elections from across the world.
For Scottish City Council and New South Wales (NSW) Local Government elections, we randomly selected one profile to audit for each of $m=2$, $3$, $4$, and $5$-seat elections. 
Elsewhere, we aimed to represent multiple scales of STV elections, and we prioritized recent elections in municipalities that still use STV today.

We recorded, for each of those elections, the maximal Margin of Insecurity $M$ for which a coherent and plausible graph could be constructed, and the Average Sampling Number (ASN) required for a successful audit.
We artificially noised these election CVRs with $2\%$ mismatches, following a standard noising procedure established in the Appendix of \cite{blomRAIRERiskLimitingAudits2019b}. We used a risk level of $\alpha = 5\%$ for each of these audits.

For the Mismatch method, which has arbitrary stopping times, we computed the ASN as the average sample size of 10 audits of the election, assuming all 10 certified with a sample size of less than $N/2$. 
For the Delta method, we defined the ASN as the smallest sample size which would lead to a successful audit in at least 9 out of 10 randomly seeded trials.
\begin{table}[h]
\centering
\caption{Margins and Sample Sizes for Graph-Based Audits of STV Elections.}
\label{tab:results}

\resizebox{\textwidth}{!}{%
\begin{tabular}{@{}clrrrrrrr@{}}
\hline
Category & Election & Year & $C$ & $m$ & $N$
& $M$
& Mismatch ASN
& Delta ASN \\
\hline
\hline

\multirow[c]{4}{*}{\makecell[c]{Scottish\\Local\\Elections}}
& Sgire Nan Loch
& 2022 & 4 & 2 & 851
& 141 (17\%)
& 66 (7.8\%)
& 56 (6.6\%) \\
\cline{2-9}

& Bearsden North
& 2017 & 6 & 3 & 7,035
& 87 (1.2\%)
& X
& 2,345 (33.3\%) \\
\cline{2-9}

& Ardrossan and Arran
& 2012 & 10 & 4 & 5,590
& 75 (1.3\%)
& X
& 1,118 (20\%) \\
\cline{2-9}

& Saltcoats and Stevenston
& 2022 & 8 & 5 & 6,506
& 148 (2.3\%)
& X
& 650 (10\%) \\
\hline\hline

\multirow[c]{4}{*}{\makecell[c]{NSW\\Local\\Elections}}
& Shellharbour Ward D
& 2024 & 6 & 2 & 11,600
& 1,933 (17\%)
& 51.7 (0.45\%)
& 46 (0.4\%) \\
\cline{2-9}

& Cumberland Greystanes
& 2017 & 13 & 3 & 19,422
& 21 (0.11\%)
& X
& X \\
\cline{2-9}

& Wollongong Ward 3
& 2021 & 16 & 4 & 40,706
& 2,200\textsuperscript{2} (5\%)
& X
& 1,017 (2.5\%) \\
\cline{2-9}

& Penrith South
& 2024 & 15 & 5 & 37,778
& 600 (1.6\%)
& X
& 3,777 (10\%) \\
\hline\hline

\multirow[c]{2}{*}{Albany, CA}
& City Council
& 2022 & 5 & 2 & 7,159
& 129 (2\%)
& X
& 715 (10\%) \\
\cline{2-9}

& City Council
& 2024 & 5 & 2 & 8,007
& 192 (2\%)
& X
& 533 (6.7\%) \\
\hline\hline



\multirow[c]{4}{*}{\makecell[c]{Portland, OR\\City Council}}
& District 1
& 2024 & 16 & 3 & 42,678
& 671 (1.6\%)
& X
& 711 (1.7\%) \\
\cline{2-9}

& District 2
& 2024 & 22 & 3 & 77,024
& 2,739\textsuperscript{1} (3.6\%)
& 581 (0.75\%)
& 308 (0.4\%) \\
\cline{2-9}

& District 3
& 2024 & 30 & 3 & 84,356
& 300 (0.36\%)
& X
& 4,217 (5\%) \\
\cline{2-9}

& District 4
& 2024 & 30 & 3 & 76,566
& 200 (0.26\%)
& X
& 7,656 (10\%) \\
\hline\hline

\multirow[c]{4}{*}{Minneapolis, MN}
& Parks Board
& 2025 & 8 & 3 & 113,348
& 3,209 (2.8\%)
& X
& 453 (0.4\%) \\
\cline{2-9}

& Parks Board
& 2021 & 7 & 3 & 106,650
& 1,111 (1\%)
& X
& 1,777 (1.7\%) \\
\cline{2-9}

& Board of Taxation
& 2025 & 3 & 2 & 106,682
& 8,467 (7.9\%)
& 425.3 (0.4\%)
& 133 (0.12\%) \\
\cline{2-9}

& Board of Taxation
& 2021 & 4 & 2 & 95,625
& 5,406 (5.7\%)
& 1,528.4 (1.6\%)
& 159 (0.17\%) \\
\hline\hline

\multirow[c]{8}{*}{\makecell[c]{Australian\\Senate}}
& ACT
& 2025 & 14 & 2 & 293,474
& 15,000 (5\%)
& 348.3 (0.1\%)
& 293 (0.1\%) \\
\cline{2-9}

& Northern Territory
& 2025 & 17 & 2 & 106,807
& 1,500 (1\%)
& 1,932.8 (2\%)
& 1,068 (1\%) \\
\cline{2-9}

& Victoria
& 2025 & 65 & 6 & 4,101,762
& 20,000\textsuperscript{1} (0.5\%)
& X
& 20,508 (0.5\%) \\
\cline{2-9}

& New South Wales
& 2025 & 56 & 6 & 4,986,832 & 12,000\textsuperscript{1} (0.2\%)
& X & 49,868 (1\%) \\
\cline{2-9}

& Tasmania
& 2025 & 33 & 6 & 371,790
& 5,000\textsuperscript{1} (1\%)
& 3,253.9 (1\%)
& 1,487 (0.4\%) \\
\cline{2-9}

& South Australia
& 2025 & 40 & 6 & 1,164,072 &
7,000\textsuperscript{1} (1\%) & X & 11,640 (1\%) \\
\cline{2-9}

& Queensland
& 2025 & 56 & 6 & 3,224,436 & 30,000\textsuperscript{1} (1\%)
& X & 8,061 (0.2\%) \\
\cline{2-9}

& Western Australia
& 2025 & 49 & 6 & 1,622,016 & 2,850\textsuperscript{1} (0.2\%)
& X & 108,134 (7\%) \\
\hline

\end{tabular}%
}

\vspace{2pt}

\begin{minipage}{\textwidth}
\footnotesize
\textsuperscript{1} The audit graph was constructed with a batch elimination.\quad
\textsuperscript{2} For Wollongong Ward 3, \(M=500\) was used for audits.
\end{minipage}
\vspace{-12pt}
\end{table}

There are two potential failure points for graph-based STV audits -- namely graph construction and assertion testing -- and the data provides insight into both.

When graph construction fails, it is usually because of combinatorial explosion of the early layers of the graph for elections with $C\geq 20$ candidates. 
The batch elimination shortcut of \S\ref{sec:shortcut} works to circumvent this problem in some cases, but it does not apply in situations where a strong candidate is close to making a quota on initial preferences, as in Wollongong or Portland District 3.
The latter is factually one of the most stable in Portland -- the three winners are close to making a solid coalition -- but because batch elimination does not apply, the only remedy is to construct a full plausible graph with a much smaller margin $M$.

It is generally possible to construct a plausible graph by pushing $M$ low enough, at the cost of a higher ASN for the audit it generates.
This tradeoff is best illustrated by the Cumberland Greystanes election, which was simply a close race: the last round was decided by a margin of 21 out of the nearly 20,000 votes cast.
This closeness was responsible for the only outright failure of the Delta method.

The Mismatch Audit ASN was only comparable to the Delta-method Audit in profiles where the diluted margin $M/N$ was large, as is predicted by the theory of the COBRA test process in \S 3 of \cite{spertusCOBRAComparisonOptimalBetting2023}.
In fairness to the Mismatch Audit, its success rates are much higher when noise levels are closer to the real-world measured rates (typically $<0.2\%$).
We chose a deliberately higher noise level so that our results might demonstrate reliability as well as effectiveness.

These results indicate how STV RLAs might further be improved -- e.g. devising more shortcuts to avoid full-graph computation -- but they are also cause for optimism. In particular, this is the first published RLA of an Australian 6-seat Senate election, the largest use of STV in the world.

Replication code for Table \ref{tab:results} is available at.
